% ==========================================================================
% Chapter 5 — HRL Architecture: The "Meta-Controller" Approach
% ==========================================================================
\section{Hierarchical Reinforcement Learning Architecture}
\label{ch:hrl_architecture}

This chapter details the second integration strategy: a Hierarchical
Reinforcement Learning (HRL) framework in which the LLM acts as a
\emph{Meta-Controller}, selecting discrete high-level \emph{Options}
that govern the behaviour of the low-level MAPPO policy.

Unlike the Mixing Board approach (\cref{ch:dynamic_implementation}),
which adjusts continuous weights over a fixed greedy planner, the
Meta-Controller replaces the planner entirely.
The LLM issues explicit macro-action assignments per agent, and the
low-level policy receives the active option as part of its observation,
creating a two-level decision hierarchy.

The exposition follows the implementation order:
\textbf{Option definitions} (\cref{sec:option_defs}) $\rightarrow$
\textbf{discrete option selection} (\cref{sec:discrete_selection}) $\rightarrow$
\textbf{one-hot observation augmentation} (\cref{sec:onehot_aug}) $\rightarrow$
\textbf{bipartite reward structure} (\cref{sec:bipartite_reward}) $\rightarrow$
\textbf{staleness and reflection} (\cref{sec:staleness}).

% ──────────────────────────────────────────────────────────────────────────
\subsection{System Overview}
\label{sec:hrl_overview}

\Cref{fig:hrl_system_flow} shows the two-level architecture.
The Meta-Controller (LLM) observes inventory transitions and
assigns Options to each agent asynchronously; the low-level MAPPO
policy conditions on the active option and receives a composite reward
signal comprising both the environment's sparse reward and
an option-specific dense intrinsic reward.

\begin{figure}[H]
\centering
\begin{tikzpicture}[
    box/.style={draw, rounded corners, minimum width=2.6cm, minimum height=0.9cm,
                font=\small\sffamily, align=center, fill=gray!8},
    llmbox/.style={draw, rounded corners, minimum width=2.6cm, minimum height=0.9cm,
                font=\small\sffamily, align=center, fill=blue!8},
    arr/.style={-{Stealth[length=2.5mm]}, thick},
    node distance=1.6cm and 2cm
]
    % Top row: LLM level
    \node[llmbox] (llm) {LLM\\Meta-Controller};
    \node[box, right=3cm of llm] (ctrl) {OptionController\\$o_t^{(0)}, o_t^{(1)}$};

    % Bottom row: MAPPO level
    \node[box, below=2cm of llm] (env) {Env Step\\$s_t \to s_{t+1}$};
    \node[box, right=3cm of env] (agent) {MAPPO Agent\\$\pi_\theta(a \mid s, o)$};

    % LLM -> Controller
    \draw[arr] (llm) -- node[above,font=\footnotesize]{JSON} (ctrl);

    % Controller -> Agent (one-hot)
    \draw[arr] (ctrl.south) -- ++(0,-0.6) -|
        node[near start,right,font=\footnotesize]{one-hot $\mathbf{e}_{o}$} (agent.north);

    % Agent -> Env
    \draw[arr] (agent.south) -- ++(0,-0.4) -|
        node[near start,below,font=\footnotesize]{$a_t$} (env.south);

    % Env -> Agent (reward)
    \draw[arr] (env.east) -- node[below,font=\footnotesize]{$r_t^{\text{env}} + r_t^{\text{intr}}$} (agent.west);

    % Env -> LLM (trigger)
    \draw[arr,dashed] (env.north) -- node[left,font=\footnotesize]{\scriptsize success / stale} (llm.south);

    % Env -> Controller (reset on terminal)
    \draw[arr,dashed] (env.north east) -- node[right,font=\footnotesize,pos=0.3]{\scriptsize terminal} (ctrl.south west);
\end{tikzpicture}
\caption{Two-level HRL architecture.
         The LLM Meta-Controller is triggered by option success or
         staleness (dashed). The OptionController passes one-hot
         embeddings to the MAPPO agent at every step (solid).}
\label{fig:hrl_system_flow}
\end{figure}


% ──────────────────────────────────────────────────────────────────────────
\subsection{Option Definitions}
\label{sec:option_defs}

An \emph{Option} $o\in\mathcal{O}$ is a macro-action defined by the
triple $\langle \mathcal{I}_o, \pi_o, \beta_o \rangle$: an initiation
set, an intra-option policy, and a termination condition~\cite{Sutton_Precup_Singh_1999}.
In our implementation, the intra-option policy is not explicitly
separate; instead, the low-level MAPPO policy is conditioned on the
active option via observation augmentation (\cref{sec:onehot_aug}).

The option set is hard-coded from the DAG:
\begin{equation}
\label{eq:option_set}
\mathcal{O} = \left\{
\begin{array}{l}
    \texttt{COLLECT\_WOOD},\;
    \texttt{COLLECT\_STONE},\;
    \texttt{CRAFT\_PICKAXE},\\
    \texttt{MINE\_IRON},\;
    \texttt{CRAFT\_SWORD},\;
    \texttt{CRAFT\_ARMOR},\\
    \texttt{BUILD\_BRIDGE},\;
    \texttt{DEFEAT\_ENEMY},\;
    \texttt{MINE\_GOLD},\;
    \texttt{IDLE}
\end{array}
\right\}
\end{equation}
giving $|\mathcal{O}|=10$ discrete options.

Each option has three associated properties defined in \texttt{options.py}:

\begin{table}[H]
\centering
\caption{Option properties: target zone, responsible agent, and initiation condition.}
\label{tab:option_props}
\small
\begin{tabular}{@{}lccc@{}}
\toprule
\textbf{Option} & \textbf{Zone} & \textbf{Agent} & \textbf{Initiation Condition} \\
\midrule
\texttt{COLLECT\_WOOD}  & Wood (0)      & 0     & Always \\
\texttt{COLLECT\_STONE} & Stone (1)     & 1     & Always \\
\texttt{CRAFT\_PICKAXE} & Workbench (2) & Both  & wood $\ge 1$, stone $\ge 1$, pickaxe $< 1$ \\
\texttt{MINE\_IRON}     & Iron (3)      & 1     & pickaxe $\ge 1$ \\
\texttt{CRAFT\_SWORD}   & Workbench (2) & Both  & iron $\ge 1$, sword $< 1$ \\
\texttt{CRAFT\_ARMOR}   & Workbench (2) & Both  & iron $\ge 1$, armor $< 1$ \\
\texttt{BUILD\_BRIDGE}  & Bridge (4)    & Both  & wood $\ge 1$, bridge $< 1$ \\
\texttt{DEFEAT\_ENEMY}  & Enemy (5)     & Both  & bridge, sword, armor $\ge 1$ \\
\texttt{MINE\_GOLD}     & Gold (6)      & 1     & enemy defeated, gold $< 1$ \\
\texttt{IDLE}           & ---           & Both  & Always \\
\bottomrule
\end{tabular}
\end{table}

\noindent
\textbf{Termination.}
An option terminates when the relevant inventory slot increases:
\begin{lstlisting}[style=code, language=Python, caption={Termination check (excerpt from \texttt{options.py}).}, label={lst:option_term}]
def is_terminated(option, inventory):
    if option == Option.COLLECT_WOOD:
        return inventory[:, I_WOOD] >= 1
    elif option == Option.CRAFT_PICKAXE:
        return inventory[:, I_PICKAXE] >= 1
    # ... analogous for all other options
    elif option == Option.IDLE:
        return np.zeros(n, dtype=bool)  # never terminates
\end{lstlisting}


% ──────────────────────────────────────────────────────────────────────────
\subsection{Discrete Option Selection by the LLM}
\label{sec:discrete_selection}

The Meta-Controller is an LLM (Qwen~2.5 7B via Ollama, or a LoRA-finetuned variant)
that receives a structured prompt and returns a JSON object specifying
one option per agent.

\subsubsection{Trigger Conditions}

Option selection is triggered under two conditions, evaluated per-environment:
\begin{enumerate}[nosep]
    \item \textbf{Success}: at least one agent's active option has completed
          (the corresponding inventory slot changed).
    \item \textbf{Staleness}: the active option has been running for more
          than $T_{\text{stale}}=150$ steps without completion.
\end{enumerate}

Environments currently awaiting an LLM response (\texttt{llm\_pending})
or on cooldown (50~steps after the last query) are excluded:
\begin{lstlisting}[style=code, language=Python, caption={Trigger logic in \texttt{hrl\_train\_loop.py}.}, label={lst:hrl_trigger}]
success_mask = a0_success | a1_success
stale_mask = option_controller.get_stale_mask(threshold=150)
ready_mask = (success_mask | stale_mask) \
             & (~option_controller.llm_pending) \
             & (option_controller.cooldown_counter == 0)
trigger_envs = np.where(ready_mask)[0]
\end{lstlisting}

\subsubsection{Prompt Structure}

For a normal trigger (success), the \texttt{PromptBuilder} constructs a
Chain-of-Thought prompt:
\begin{lstlisting}[style=code, language=Python, caption={HRL prompt schema (from \texttt{prompt\_builder.py}).}, label={lst:hrl_prompt}]
"""
You are the high-level Cognitive Orchestrator for two MARL agents.
Strict Dependency DAG: (Wood+Stone) -> Pickaxe -> Iron
    -> (Sword+Armor) -> Bridge -> Enemy -> Gold.

### CURRENT STATE:
Inventory: {inventory}
Agent 0 Status: {a0_status}
Agent 1 Status: {a1_status}

### YOUR TASK:
Assign exactly ONE discrete Option to each agent.
Available Options: ["COLLECT_WOOD", "COLLECT_STONE", ...]

Respond ONLY with valid JSON:
{
  "dag_check": "<1 sentence reasoning>",
  "agent_0_option": "<Option>",
  "agent_1_option": "<Option>"
}
"""
\end{lstlisting}

For a staleness trigger, a \emph{counterfactual reflection prompt} is
used instead, forcing the LLM to diagnose \textbf{why} the previous
assignment failed before issuing new options:
\begin{lstlisting}[style=code, language=Python, caption={Reflection prompt for stale options (from \texttt{prompt\_builder.py}).}, label={lst:reflection_prompt}]
"""
### REFLECTION REQUIRED:
Agent 0 was assigned "{a0_option}" and Agent 1
was assigned "{a1_option}".
After {stale_steps} steps, NEITHER has succeeded.

### YOUR TASK:
1. First, diagnose what prerequisite is MISSING.
2. Then assign corrected options that respect the DAG.

Respond ONLY with JSON:
{
  "reflection": "<diagnosis>",
  "dag_check": "<verification>",
  "agent_0_option": "<Option>",
  "agent_1_option": "<Option>"
}
"""
\end{lstlisting}

\subsubsection{Batched Queries with Deduplication}

Multiple environments may trigger simultaneously. To avoid redundant
LLM calls, environments with identical inventory states are
deduplicated: the system groups environments by unique inventory
vector and sends one prompt per unique state. The response is then
broadcast to all environments sharing that state:

\begin{lstlisting}[style=code, language=Python, caption={Inventory deduplication in \texttt{hrl\_train\_loop.py}.}, label={lst:dedup}]
inv_batch = inv_next[trigger_envs].astype(int)
unique_invs, inverse_indices = np.unique(
    inv_batch, axis=0, return_inverse=True
)
unique_to_envs = [
    trigger_envs[inverse_indices == i]
    for i in range(len(unique_invs))
]
\end{lstlisting}

All unique prompts are dispatched as a batch to the \texttt{LLMBridge}
via \texttt{query\_batch\_async}, and responses are applied asynchronously.

\subsubsection{Environment Freezing During Pending Queries}

While an environment is awaiting an LLM response, it is frozen to
prevent state drift:
\begin{lstlisting}[style=code, language=Python, caption={Environment freezing for pending queries.}, label={lst:env_freeze}]
pending_mask = option_controller.llm_pending.copy()
if pending_mask.any():
    actions_np[pending_mask] = 4  # NO-OP
    saved_pos = vec_env.pos[pending_mask].copy()
    saved_inv = vec_env.inventory[pending_mask].copy()
    # ... after env.step(), restore saved state:
    vec_env.pos[pending_mask] = saved_pos
    vec_env.inventory[pending_mask] = saved_inv
    vec_env.step_counts[pending_mask] -= 1  # pause timer
\end{lstlisting}

This ensures that the LLM's decision is applied to the exact state
it observed, not a state that has drifted during inference latency.

% ──────────────────────────────────────────────────────────────────────────
\subsection{One-Hot Observation Augmentation}
\label{sec:onehot_aug}

The active option is communicated to the low-level MAPPO policy via
a one-hot vector appended to the agent's observation.

\subsubsection{Observation Construction}

The vector input to the MAPPO network in HRL mode is:
\begin{equation}
\label{eq:hrl_obs}
    \mathbf{v}_t^{(i)}
    =
    \left[
        \underbrace{\mathbf{p}^{(i)}}_{\mathbb{R}^2}
        \;\Big\|\;
        \underbrace{\mathbf{0}}_{\mathbb{R}^3\text{ (pad)}}
        \;\Big\|\;
        \underbrace{\mathbf{x}}_{\mathbb{R}^{10}\text{ (inv)}}
        \;\Big\|\;
        \underbrace{\mathbf{e}_{o_t^{(i)}}}_{\mathbb{R}^{10}\text{ (one-hot)}}
    \right]
    \;\in\; \mathbb{R}^{25}
\end{equation}
where $\mathbf{p}^{(i)}$ is the agent's position,
the zero-padding replaces the goal embedding from the Mixing Board
approach (maintaining dimensional compatibility),
$\mathbf{x}$ is the inventory vector, and
$\mathbf{e}_{o}\in\{0,1\}^{|\mathcal{O}|}$ is the one-hot encoding
of the active option.

\begin{lstlisting}[style=code, language=Python, caption={Observation construction in \texttt{hrl\_train\_loop.py}.}, label={lst:hrl_obs}]
goal_emb = np.zeros((n_envs, 2, 3), dtype=np.float32)
inv_repeat = np.stack([
    all_obs[:, 0, 4:4+NUM_ITEMS],
    all_obs[:, 0, 4:4+NUM_ITEMS]
], axis=1)
opt_repeat = option_controller.get_option_embeddings()
pos_repeat = vec_env.pos.copy()

vec_input = np.concatenate(
    [pos_repeat, goal_emb, inv_repeat, opt_repeat], axis=2
)
\end{lstlisting}

\subsubsection{One-Hot Embedding Generation}

The \texttt{OptionController} maintains per-environment, per-agent
option indices and produces embeddings via vectorised indexing:

\begin{lstlisting}[style=code, language=Python, caption={One-hot embedding in \texttt{option\_controller.py}.}, label={lst:onehot_gen}]
def get_option_embeddings(self):
    """Returns [n_envs, 2, NUM_OPTIONS] one-hot."""
    embs = np.zeros((self.n_envs, 2, NUM_OPTIONS),
                    dtype=np.float32)
    env_idx = np.arange(self.n_envs)
    embs[env_idx, 0, self._active_options_a0] = 1.0
    embs[env_idx, 1, self._active_options_a1] = 1.0
    return embs
\end{lstlisting}

This one-hot conditioning allows the same MAPPO network to learn
\emph{option-conditional behaviour} without separate policy networks
per option. The policy $\pi_\theta(a \mid s, o)$ implicitly specialises
its behaviour depending on which option is active.

% ──────────────────────────────────────────────────────────────────────────
\subsection{Bipartite Reward Structure}
\label{sec:bipartite_reward}

The HRL approach uses a composite reward with two distinct components:
a \textbf{sparse} environment reward and a \textbf{dense} intrinsic
reward. This bipartite structure enables learning from both the
environment's ground-truth signal and the Meta-Controller's guidance.

\subsubsection{Environment Reward (Sparse)}
The environment reward $r_t^{\text{env}}$ is the unchanged signal from
\texttt{BatchCraftingEnvV2}: a step penalty of $-0.01$ plus milestone
bonuses upon successful interactions:

\begin{table}[H]
\centering
\caption{Environment reward structure from \texttt{crafting\_env.py}.}
\label{tab:env_rewards}
\begin{tabular}{@{}lc@{}}
\toprule
\textbf{Event} & \textbf{Reward (shared)} \\
\midrule
Step penalty      & $-0.01$ \\
Collect Wood      & $+2.0$  \\
Collect Stone     & $+2.0$  \\
Craft Pickaxe     & $+3.0$  \\
Mine Iron         & $+2.0$  \\
Craft Sword/Armor & $+3.0$  \\
Build Bridge      & $+3.0$  \\
Defeat Enemy      & $+10.0$ \\
Mine Gold (win)   & $+15.0$ \\
\bottomrule
\end{tabular}
\end{table}

\subsubsection{Intrinsic Reward (Dense)}
\label{sec:intrinsic_reward}

The intrinsic reward provides a dense shaping signal that guides agents
toward the target zone of their currently active option. It consists of
two components:

\paragraph{Distance-based shaping.}
For each agent, the potential-based shaping reward toward the active
option's target zone is computed identically to the PBRS
in~\cref{sec:pbrs}, but applied to the option's target rather than
the planner's target:
\begin{equation}
\label{eq:hrl_shaping}
    r_t^{\text{shape},(i)}
    =
    \begin{cases}
        \bigl(\gamma\,\Phi^{(i)}(s_{t+1}, z_o) - \Phi^{(i)}(s_t, z_o)\bigr)
            \cdot \lambda_{\text{intr}}
        & \text{if } o^{(i)} \neq \texttt{IDLE} \\[4pt]
        0
        & \text{if } o^{(i)} = \texttt{IDLE}
    \end{cases}
\end{equation}
where $z_o$ is the target zone of option~$o$
(\cref{tab:option_props}), $\Phi$ is the same potential function from
\cref{eq:potential}, and $\lambda_{\text{intr}} = 0.05$.

The vectorised implementation:
\begin{lstlisting}[style=code, language=Python, caption={Per-agent intrinsic shaping in \texttt{hrl\_train\_loop.py}.}, label={lst:hrl_intrinsic}]
z0 = get_option_target_zone(a0_opt)
idle_mask_0 = (z0 == -1)
safe_z0 = np.where(idle_mask_0, 0, z0)

dist0_prev = np.linalg.norm(
    pos_prev[:, 0] - vec_env.zones[np.arange(n_envs), safe_z0],
    axis=1)
dist0_next = np.linalg.norm(
    vec_env.pos[:, 0] - vec_env.zones[np.arange(n_envs), safe_z0],
    axis=1)
phi0_prev = MAX_DIST - dist0_prev
phi0_next = MAX_DIST - dist0_next
shaping0 = np.where(
    idle_mask_0, 0.0,
    (0.99 * phi0_next) - phi0_prev
)
intrinsic_r[:, 0] = np.where(
    a0_success, 1.0, shaping0 * 0.05
)
\end{lstlisting}

\paragraph{Option success bonus.}
When an agent completes its active option (the relevant inventory slot
increases), it receives a bonus of $+1.0$:
\begin{equation}
\label{eq:success_bonus}
    r_t^{\text{intr},(i)}
    =
    \begin{cases}
        1.0 & \text{if option } o^{(i)} \text{ terminated successfully} \\[2pt]
        r_t^{\text{shape},(i)} & \text{otherwise}
    \end{cases}
\end{equation}

Success is detected by comparing the inventory before and after the
environment step:
\begin{lstlisting}[style=code, language=Python, caption={Option success detection in \texttt{hrl\_train\_loop.py}.}, label={lst:success_check}]
def check_option_success(opt_strs, inv_prev, inv_next):
    for i in range(n_envs):
        opt = opt_strs[i]
        if opt == "COLLECT_WOOD":
            success[i] = inv_next[i, 0] > inv_prev[i, 0]
        elif opt == "CRAFT_PICKAXE":
            success[i] = inv_next[i, 3] > inv_prev[i, 3]
        # ... analogous for all options
    return success
\end{lstlisting}

\subsubsection{Composite Reward}

The total reward signal is a weighted sum of the environment and
intrinsic components:
\begin{equation}
\label{eq:hrl_total_reward}
\boxed{
    r_t^{\text{total}}
    = r_t^{\text{env}}
    + \alpha_{\text{intr}} \cdot r_t^{\text{intr}}
}
\end{equation}
with $\alpha_{\text{intr}} = 0.5$. This coefficient was chosen to
balance the dense intrinsic signal against the sparse environment
milestones; too large a value causes the policy to over-optimise
for intrinsic shaping at the expense of actual task completion.

\begin{lstlisting}[style=code, language=Python, caption={Composite reward in \texttt{hrl\_train\_loop.py}.}, label={lst:composite_reward}]
total_r = env_rewards + 0.5 * intrinsic_r
\end{lstlisting}

Note the key difference from the Mixing Board approach: here, no
$\lambda$ decay schedule is applied. The intrinsic coefficient
remains fixed because the LLM continuously updates the \emph{direction}
(which option to pursue) rather than the \emph{magnitude} of shaping.

% ──────────────────────────────────────────────────────────────────────────
\subsection{Staleness Detection and Counterfactual Reflection}
\label{sec:staleness}

\subsubsection{Staleness Detection}

The \texttt{OptionController} tracks the \emph{age} (number of steps)
of each active option per environment. If an option exceeds
$T_{\text{stale}}=150$ steps without termination, the environment
is marked as stale and the LLM is re-queried:

\begin{lstlisting}[style=code, language=Python, caption={Staleness detection in \texttt{option\_controller.py}.}, label={lst:stale_detect}]
def get_stale_mask(self, threshold=150):
    IDLE_IDX = OPTION_NAMES.index("IDLE")
    age_exceeded = self.option_age >= threshold
    not_idle = (self._active_options_a0 != IDLE_IDX) \
             | (self._active_options_a1 != IDLE_IDX)
    return age_exceeded & not_idle
\end{lstlisting}

Crucially, \texttt{IDLE} options are excluded from staleness triggers,
because intentional waiting (e.g.\ Agent~0 waiting while Agent~1
collects stone) is a valid strategy.

\subsubsection{Counterfactual Reflection}

When a staleness trigger fires, the LLM receives a reflection prompt
instead of the standard assignment prompt. This forces the LLM to
reason about \emph{why} the previous option failed—typically a missing
DAG precondition—before proposing a corrected assignment.

This design addresses the observed failure mode where the LLM
repeatedly assigned options whose preconditions were not met (e.g.\
\texttt{MINE\_IRON} without a pickaxe). The reflection mechanism
leverages the LLM's zero-shot reasoning to debug its own prior
hallucination.

% ──────────────────────────────────────────────────────────────────────────
\subsection{Episode Lifecycle and Option Reset}
\label{sec:episode_lifecycle}

When an environment reaches a terminal state (gold mined, game over,
or truncation at 800 steps), the \texttt{OptionController} resets
that environment to its initial options:

\begin{lstlisting}[style=code, language=Python, caption={Episode reset in \texttt{option\_controller.py}.}, label={lst:option_reset}]
def reset_options(self, env_indices):
    self._active_options_a0[env_indices] = 0  # COLLECT_WOOD
    self._active_options_a1[env_indices] = 1  # COLLECT_STONE
    self.llm_pending[env_indices] = False
    self.cooldown_counter[env_indices] = 0
    self.option_age[env_indices] = 0
\end{lstlisting}

The initial options (\texttt{COLLECT\_WOOD} for Agent~0,
\texttt{COLLECT\_STONE} for Agent~1) correspond to the first layer
of the DAG, ensuring agents begin collecting resources immediately.

% ──────────────────────────────────────────────────────────────────────────
\subsection{PPO Integration: Target KL Guard}
\label{sec:target_kl}

The HRL training loop adds a \emph{target KL divergence} early-stopping
mechanism to PPO. During each PPO epoch, the approximate KL divergence
between the old and new policies is monitored:
\begin{equation}
\label{eq:approx_kl}
    D_{\text{KL}}^{\text{approx}}
    = \mathbb{E}\bigl[-\log \pi_{\theta_{\text{new}}}(a|s)
      + \log \pi_{\theta_{\text{old}}}(a|s)\bigr]
\end{equation}
If $D_{\text{KL}}^{\text{approx}} > 0.015$, the remaining PPO
epochs are aborted for that update:

\begin{lstlisting}[style=code, language=Python, caption={Target KL guard in \texttt{hrl\_train\_loop.py}.}, label={lst:target_kl}]
with torch.no_grad():
    approx_kl = (-logratio).mean().item()
target_kl = 0.015
if approx_kl > target_kl:
    target_kl_reached = True
    break
\end{lstlisting}

This guard prevents destructive policy updates when the intrinsic
reward signal changes abruptly after an LLM option reassignment.

% ──────────────────────────────────────────────────────────────────────────
\subsection{Comparison with the Mixing Board Approach}
\label{sec:hrl_vs_dyn}

\begin{table}[H]
\centering
\caption{Structural comparison of the two LLM integration strategies.}
\label{tab:approach_comparison}
\small
\begin{tabular}{@{}p{3.5cm}p{5cm}p{5cm}@{}}
\toprule
\textbf{Aspect} & \textbf{Mixing Board (Ch.~4)} & \textbf{Meta-Controller (Ch.~5)} \\
\midrule
LLM output space &
    Continuous logits $\in \mathbb{R}^7$ &
    Discrete options $\in \mathcal{O}^2$ \\
LLM controls &
    Subtask priority weights &
    Explicit per-agent macro-actions \\
Trigger mechanism &
    Plateau + Z-score on TD error &
    Option success or staleness \\
Query granularity &
    Global (once per trigger) &
    Per-environment (batched) \\
Observation augmentation &
    Goal embedding $\in \mathbb{R}^3$ &
    One-hot option $\in \{0,1\}^{10}$ \\
Reward structure &
    PBRS with $\lambda$ decay &
    Bipartite: sparse env + dense intrinsic (fixed $\alpha$) \\
LLM query frequency &
    Low (plateau-triggered) &
    High (per-option completion) \\
Safety mechanism &
    Softmax budget, EMA, rollback &
    DAG in prompt, reflection, env freezing \\
\bottomrule
\end{tabular}
\end{table}


% ──────────────────────────────────────────────────────────────────────────
\subsection{Summary}
\label{sec:hrl_summary}

The Meta-Controller approach creates a clean separation of concerns:
the LLM reasons about \emph{what to do} (strategic option selection),
while the MAPPO policy learns \emph{how to do it} (low-level
navigation and interaction). The one-hot conditioning allows a single
network to exhibit option-dependent behaviour, and the bipartite
reward structure ensures that the policy receives both a consistent
dense gradient from intrinsic shaping and a grounding signal from
the true environment rewards.

\begin{table}[H]
\centering
\caption{Key hyperparameters of the Meta-Controller approach.}
\label{tab:hrl_hyperparams}
\begin{tabular}{@{}lll@{}}
\toprule
\textbf{Component} & \textbf{Parameter} & \textbf{Value} \\
\midrule
Options          & $|\mathcal{O}|$                & 10 \\
Staleness        & $T_{\text{stale}}$             & 150 steps \\
Cooldown         & Post-query cooldown            & 50 steps \\
Intrinsic reward & Distance scale $\lambda_{\text{intr}}$ & 0.05 \\
                 & Success bonus                  & $+1.0$ \\
Composite reward & $\alpha_{\text{intr}}$         & 0.5 \\
                 & $D_{\max}$                     & 85.0 \\
Target KL        & KL threshold                   & 0.015 \\
Network          & One-hot dimension              & 10 \\
                 & Total \texttt{vec\_dim}        & 25 \\
\bottomrule
\end{tabular}
\end{table}
